\section{Cote 123: a companion to the letter of 22 October 1969}
\label{sec:arch123}

\paragraph{Scope and provenance.}
The two pages of the letter are physical PDF pp.~22--23, archive
footers 21--22, supplied files \texttt{ARCHIVE\_123\_scan022.tex} and
\texttt{ARCHIVE\_123\_scan023.tex}.  Scan 024 is the short following
heading at footer 23, not a third page of the letter.  The first page
concerns a complex simple adjoint group; the qualification ``simply
connected instead of adjoint'' occurs on the second page specifically
in the question about rational singularities of the group quotient.
Throughout this section automorphisms are complex algebraic
automorphisms.  All mathematical inputs below are published theorems
or the explicitly accepted ALPOGE\_BRIDGES checkpoint~09 statements;
none uses Alp\"oge, Masullo, or the three outside CM preprints.

This section consolidates accepted results rather than presenting them
as new.  Its new checked deductions are the verification
\textbf{ARCH-123-NORMAL} and the refinement
\textbf{ARCH-123-COUNT}.  The former closes the precise normal-form
check requested for CC-APG04's converse.  Historical novelty is not
claimed for either deduction.

\subsection{The classification and rigidity questions (scan 022)}
Let $G$ be simple and adjoint over $\mathbf C$, let $r$ be its rank,
$h$ its Coxeter number, $W$ its Weyl group, and
$z=|Z(G_{\mathrm{sc}})|$.  An ordered pair $(T_0,T_1)$ is in
principal apposition if there is a principal element $c\in T_0$ of
order $h$ such that
\[
 Z_G(c)^\circ=T_0,\qquad c\in N_G(T_1),\qquad
 C=\operatorname{Ad}(c)|_{\operatorname{Lie}T_1}
 \text{ is a Coxeter element}.
\]
The identity component is essential: the full centralizer of a
regular semisimple element in an adjoint group need not be connected.
The convention is Kostant (1959), Theorem~8.6 and the immediately
following Corollary~8.6, p.~1026; a directly available account is
Kostant (2010), Theorems~1.2--1.3, p.~6, the principal-element paragraph
on p.~7, and Proposition~1.8--Theorem~1.9, p.~8.

For the fixed ordered pair put
\[
\begin{split}
 W_i&=N_G(T_i)/T_i,\\
 U_h&=(\mathbf Z/h\mathbf Z)^\times,\\
 A_h&=\{w\in W_0:w(c)=c^a\text{ for some }a\in U_h\},\\
 N_h&=N_{W_1}(\langle C\rangle).
\end{split}
\]
The two multiplier maps $\mu_0:A_h\to U_h$ and
$\mu_1:N_h\to U_h$ are defined respectively by $w(c)=c^{\mu_0(w)}$
and $wCw^{-1}=C^{\mu_1(w)}$.  In integral root-datum coordinates,
when $c=\exp(2\pi i\rho^\vee/h)$,
\[
 A_h=\{w:w\rho^\vee-a\rho^\vee\in hP^\vee
                  \text{ for some }a\in U_h\},
 \qquad P^\vee=X_*(T_0).
\]
This makes the following classification an effective finite one.

\paragraph{Accepted result: CC-APG01--02.}
All ordered apposition pairs form one $G$-orbit,
$T_0\cap T_1=1$, and
\[
 K=N_G(T_0)\cap N_G(T_1)
   \simeq A_h\times_{U_h}N_h.
\]
There is an exact sequence
\[
 1\longrightarrow \langle c\rangle\times (T_1)^c
 \longrightarrow K\longrightarrow U_h\longrightarrow1,
 \qquad (T_1)^c\simeq Z(G_{\mathrm{sc}}).
\]
The last identification is an isomorphism of finite diagonalizable
complex groups, not a specified canonical identification; no splitting
of this extension is asserted.  With base Borels chosen for the two
tori, the inner-conjugacy classes of ordered quadruples are
\[
 K\backslash(W_0\times W_1),\qquad
 |K|=hz\varphi(h),\qquad
 \#\text{classes}=\frac{|W|^2}{hz\varphi(h)}.
\]
The action is free.  If a generating principal element of
$T_0\cap N_G(T_1)=\langle c\rangle$ is included, normalize it to $c$;
then the quotient group acting is $\langle c\rangle\times(T_1)^c$
and the count is $|W|^2/(hz)$.  These are exactly the two counts
written in the letter, including its marginal normalizer extension.
The proof is in ALPOGE\_BRIDGES, checkpoint~09,
\texttt{research/consequences\_criteria\_apposition\_general.tex},
subsections CC-APG01 and CC-APG02.  One of its known ingredients,
$N_{W_1}(\langle C\rangle)/\langle C\rangle\simeq U_h$,
is explicitly stated already in Kostant (2010), \S0.2,
equation~(0.6), p.~3; it should not be presented as a new result.

Every quadruple has trivial inner automorphism group.  Indeed, an
inner automorphism preserving $(T_i,B_i)$ is represented by an element
of $T_i$, and $T_0\cap T_1=1$.  Consequently
\[
 \operatorname{Aut}(G;T_0,B_0,T_1,B_1)
 \simeq\operatorname{Stab}_{\operatorname{Out}(G)}
       [T_0,B_0,T_1,B_1].
\]
This distinguishes the trivial stabilizer inside $G$ from the possibly
nontrivial full automorphism group asked about in the letter.

\subsection{Opposition, distinguished classes, and automorphisms (scan 022)}
Two Borels are opposite when their intersection is a maximal torus.
There is always some opposite choice for any two maximal tori, even
without apposition: this is CC-APG03, with its Bruhat argument in the
same checkpoint~09 module.  The stronger answer for apposition is:
\begin{center}
\begin{tabular}{lll}
\hline
Simple type & Every containing Borel pair opposite? & Accepted locator\\
\hline
$A_r$ & Exactly when $r+1$ is prime & CC-AP03\\
$B_r,C_r$, $r\geq2$ & Exactly when $r=2$ & CC-APG06--07\\
$D_r$, $r\geq4$ & No & CC-APG08\\
$G_2$ & Yes & BR-123-OPPG2\\
$F_4$ & No & BR-123-OPPF4\\
$E_6,E_7$ & No & CC-APG09--11\\
$E_8$ & No & BR-123-OPPE8\\
\hline
\end{tabular}
\end{center}
The rank-one identifications are $A_1=B_1=C_1$, and $B_2=C_2$.
The exceptional locators are, respectively,
\nolinkurl{research/root_g2_opposition.tex},
\nolinkurl{research/root_exceptional_opposition.tex}, and
\nolinkurl{research/consequences_criteria_exceptional_minuscule.tex}
in checkpoint~09.  These accepted inputs are not reclassified by a
new computation in this run.

In type $A_{n-1}$ this table has a particularly concrete meaning.
If $F=(\zeta^{ja})_{j,a\in\mathbf Z/n}$ is the Fourier matrix,
opposition of a coordinate flag with ordering $\pi$ and a Fourier
flag with ordering $\sigma$ says that every minor with rows
$\pi(k),\ldots,\pi(n-1)$ and columns
$\sigma(0),\ldots,\sigma(n-k-1)$ is nonzero.
For prime $n$ this follows from the Chebotarev Fourier-minor theorem
in Tao (2005), Lemma~1.3, p.~123, proof pp.~124--125.
For $n=ab$ composite, rows $0,a$ and columns $0,b$ give an all-ones
$2$ by $2$ minor.  Ordering these as the last two coordinate rows and
first two Fourier columns produces a nonopposite pair.  A direct
unipotent witness is
\[
 N=(v_0-v_b)(e_0^*-e_a^*),\quad N^2=0,\quad
 1+tN\in B\cap B',\qquad
 v_c=\sum_j\zeta^{jc}e_j.
\]
These are the precise formulas of CC-AP02--03, not an inference from
``generic position'' alone.

\paragraph{Outer-invariant existence and the finite automorphism test.}
The accepted CC-APG04--05 answer is
\[
 \text{an $\operatorname{Out}(G)$-invariant class exists}
 \quad\Longleftrightarrow\quad
 \text{the type is not }A_{2m},\quad m\geq1.
\]
This is the necessary invariance condition in Grothendieck's question
about a distinguished class.  Existence does not select a unique or
canonical class.  In allowed types choose the diagram-fixed base pair
of Borels and principal element constructed in CC-APG04.  Write
$\Gamma=\operatorname{Out}(G)$.  The group of a class represented by
$(w_0,w_1)$ is exactly
\[
 \{\gamma\in\Gamma:
 (\gamma(w_0)w_0^{-1},\gamma(w_1)w_1^{-1})
               \in A_h\times_{U_h}N_h\}.
\]
Thus it is determined by two explicit finite membership tests and
equality of their multipliers.  For all $A_{2m}$ quadruples it is
trivial.  The ambient diagram group is trivial in
$A_1,B_r,C_r,E_7,E_8,F_4,G_2$; it is $C_2$ in
$A_r$ ($r\geq2$), $D_r$ ($r\geq5$), and $E_6$, and $S_3$ in $D_4$.
The finite test, rather than the list of ambient groups alone,
determines the group of an individual quadruple.

For perspective on the positive half, the construction uses a pinning
and multiplies simple Weyl representatives orbit by orbit under
$\Gamma$.  Outside $A_{2m}$ the simple roots in each such orbit are
orthogonal, so the factors commute and the product is fixed by
$\Gamma$.  Its connected centralizer is the other torus.  A torus
inside $(G^\Gamma)^\circ$ containing that regular element supplies a
$\Gamma$-fixed regular cocharacter and hence a stable Borel.  The
orthogonality fact and the same $A_{2m}$ exception occur already in
Steinberg (1965), proof of Theorem~9.4, pp.~72--73.  That theorem
constructs rational cross-sections; it does not state the negative
quadruple assertion.  The required negative half is checked next.

\subsection{ARCH-123-NORMAL: verification of CC-AP04--05}
\label{arch123normal}
\paragraph{Statement.}
For every $n\geq2$, the ordered diagonal/Fourier apposition pair in
$G=\operatorname{PGL}_n(\mathbf C)$ has simultaneous normalizer
\[
 K_n=\{[g_{a,b,d}]:a\in(\mathbf Z/n)^\times, b,d\in\mathbf Z/n\},
 \qquad g_{a,b,d}e_j=\zeta^{ad j}e_{aj+b}.
\]
Its actions on coordinate and Fourier line labels are
$j\mapsto aj+b$ and $c\mapsto a^{-1}c+d$.  They give a free action
on pairs of complete orderings.  Inverse transpose acts on those
orderings by $(\pi,\sigma)\mapsto(\pi^{\mathrm{rev}},
-\sigma^{\mathrm{rev}})$.  For odd $n\geq3$ its action on the quotient
has no fixed class.  This verifies the normal form, the composite-$n$
case, and the precise implication on which the converse in CC-APG04
rests.

\paragraph{Full proof check and compact independent derivation.}
Let $I=\mathbf Z/n$, $Qe_j=\zeta^je_j$, and
$Ce_j=e_{j+1}$.  Then $Cv_c=\zeta^{-c}v_c$ and
$Qv_c=v_{c+1}$; $T_0$ is diagonal and $T_1=FT_0F^{-1}$.
The eigenline cycle of $Q$ on $T_1$ is a Coxeter cycle, so this is
apposition.  Conversely, conjugate the second torus of any apposition
pair to the diagonal torus.  Its principal element is represented by
a monomial matrix acting as one $n$-cycle.  After relabelling it has
$Ae_j=c_je_{j+1}$, with all $c_j\ne0$.  Choose $c$ with
$c^n=\prod_jc_j$ and set $d_{j+1}=cd_j/c_j$.  The product equation
makes this recursion consistent around the cycle; diagonal
conjugation changes $A$ to $cC$.  Its connected centralizer is the
Fourier torus.  Finally $F^2=nR$, with $Re_j=e_{-j}$, exchanges the
two torus positions by conjugation with $F^{-1}$.  Thus no ordered
apposition pair is omitted.

For an element of both normalizers choose a lift $g$ and write
\[
 ge_j=u_je_{p(j)},\qquad gv_c=w_cv_{s(c)},
 \quad u_jw_c\ne0,\quad p,s\in S_I.
\]
Comparing the $p(j)$ coordinate gives
$u_j\zeta^{jc}=w_c\zeta^{p(j)s(c)}$.  Dividing by its $c=0$
instance and then by the $j=0$ instance yields
\[
 jc=(p(j)-p(0))(s(c)-s(0))\pmod n.
\]
At $j=c=1$ the two factors multiply to $1$.  Each is therefore a
unit even when $n$ is composite.  Set $a=p(1)-p(0)$,
$b=p(0)$, $d=s(0)$.  Substitution of $c=1$ and $j=1$ gives
$p(j)=aj+b$ and $s(c)=a^{-1}c+d$ for every label.  The $c=0$
equation gives $u_j=u\zeta^{ad j}$ with a common nonzero scalar
$u$, which disappears in the projective group.  Conversely direct
substitution gives
\[
 g_{a,b,d}v_c=\zeta^{-b(a^{-1}c+d)}v_{a^{-1}c+d}.
\]
Thus these are exactly all simultaneous normalizers.  Composition is
$(a,b,d)(a',b',d')=(aa',ab'+b,a^{-1}d'+d)$, so
$K_n\simeq(I\times I)\rtimes I^\times$.  An eigenline permutation
preserving every initial segment of a complete ordering is the
identity.  If both flags are fixed, both displayed affine
permutations are identity; therefore $a=1,b=d=0$ and the action is
free.

The automorphism $\theta(g)=(g^{-1})^{\mathsf T}$ sends a flag to
its annihilator flag in reverse order.  The coordinate basis is
self-dual, whereas $v_c^{\mathsf T}v_d=n\delta_{c+d,0}$; hence the
Fourier dual labels are $-c$.  Consequently a class is fixed by
$\theta$ exactly if some $(a,b,d)$ satisfies
\[
 a\pi(i)+b=\pi(n-1-i),\qquad
 a^{-1}\sigma(i)+d=-\sigma(n-1-i) \quad(0\leq i<n).
\]
If $n$ is odd, the first affine map $R(j)=aj+b$ has the cycle
structure of reversal: it is an involution with exactly one fixed
point.  Its square implies $a^2=1$, and its fixed-point equation
$(a-1)j=-b$ has exactly one solution, so $a-1$ is a unit in $I$.
It follows from $(a-1)(a+1)=0$ that $a=-1$.  The second condition
then says the translation $S(c)=c-d$ is also reversal.  Being an
involution forces $2d=0$; oddness gives $d=0$.  Identity cannot
reverse an ordering of at least three distinct labels.  This is the
contradiction, and proves the exclusion without an assumption that
$n$ is prime.  For even $n=2m\geq4$, the lists
\[
 \pi=(0,1,\ldots,m-1,2m-1,2m-2,\ldots,m),\qquad
 \sigma=(0,1,\ldots,2m-1)
\]
satisfy the equations with $a=1,b=m,d=1$.  Triviality of the inner
stabilizer means the full automorphism group is $C_2$ in the fixed
case and $1$ otherwise.  This completes the whole argument used by
CC-AP04--05 and the converse of CC-APG04.\hfill$\square$

\subsection{ARCH-123-COUNT: a type-A refinement of the automorphism answer}
\paragraph{Statement.}
For $n=2m\geq4$, write $n=2^s u$ with $u$ odd, and let $\omega(u)$
count its distinct prime divisors ($\omega(1)=0$).  Among the
$(n!)^2/(n^2\varphi(n))$ inner-conjugacy classes of apposition
quadruples, exactly
\[
 F_n=\frac{[2^m m!]^2}{n^2\varphi(n)}\,
 \begin{cases}
 (n/2)2^{\omega(u)},&s=1,\\
 n2^{\omega(u)},&s\geq2
 \end{cases}
\]
have full automorphism group $C_2$; every other class has trivial
full automorphism group.  For odd $n\geq3$, $F_n=0$.
In particular $F_4=8$ among $18$ classes and $F_6=192$ among
$7200$ classes.  Thus even in an allowed type, the invariance
condition alone usually leaves more than one possible distinguished
class.  This is a new deduction from the checked finite normal form,
not a claim of historical priority.

\paragraph{Proof.}
Let $\Omega=S_I\times S_I$ and use its free $K_n$ action.  Each
$\theta$-fixed orbit contributes exactly $|K_n|$ pairs
$(k,\omega)$ satisfying $k\omega=\theta\omega$, one for every
$\omega$ in the orbit, because $k$ is unique by freeness.
For $k=(a,b,d)$ these equations say that
\[
 R(j)=aj+b,\qquad S(c)=-a^{-1}c-d
\]
are both conjugate in $S_n$ to reversal.  When $n=2m$, reversal has
$m$ disjoint transpositions and no fixed points.  Each such $R$
is conjugate to reversal by exactly $2^m m!$ orderings, the order
of the centralizer of this permutation.  The same count holds
independently for $S$.  Thus
\[
 |K_n|F_n=t_n[2^m m!]^2,
\]
where $t_n$ counts triples for which $R,S$ are fixed-point-free
involutions.

The involution equations are
$a^2=1$, $(a+1)b=0$, and $(a-1)d=0$ modulo $n$.
The fixed-point equations are $(a-1)j=-b$ and
$(a+1)c=-ad$.  By the Chinese remainder theorem, at each odd
prime power $p^e\Vert u$ the only solutions of $a^2=1$ are
$a=\pm1$.  For either choice the number of pairs $(b,d)$
satisfying the involution equations is $p^e$ (one variable is
arbitrary, the other zero).  An involution of a set of odd size
has a fixed point.  Therefore fixed-point-freeness of $R$ and $S$
modulo $n$ is equivalent to fixed-point-freeness of both modulo
$2^s$.  The total odd part of the count is $u2^{\omega(u)}$.

For $s=1$, $a=1$ and precisely $b=d=1$ work, giving one
triple.  For $s=2$, $a=\pm1$: when $a=1$, $b=2$ and $d$
is odd; when $a=-1$, $d=2$ and $b$ is odd.  The count is four.
For $s\geq3$, the four square roots of $1$ modulo $2^s$ are
$1,-1,1+2^{s-1},-1+2^{s-1}$.  If $a=1$, the first affine
map is a nonidentity translation of order two exactly when
$b=2^{s-1}$, and the second reflection has no fixed point
exactly when $d$ is odd: this gives $2^{s-1}$ triples.  The
case $a=-1$ gives another $2^{s-1}$ with $b$ odd and
$d=2^{s-1}$.  If $a=1+2^{s-1}$, the equation $(a+1)b=0$
forces $b$ to be a multiple of $2^{s-1}$, so the first
fixed-point equation is solvable.  If $a=-1+2^{s-1}$,
$(a-1)d=0$ forces $d$ to be a multiple of $2^{s-1}$,
so the second fixed-point equation is solvable.  These two
values contribute nothing.  Thus the $2$-part count is
$1$ for $s=1$ and $2^s$ for $s\geq2$.  Multiplying by
the odd-part count gives $t_n$ as displayed, and dividing by
$|K_n|=n^2\varphi(n)$ proves the formula.  The odd case was
proved in ARCH-123-NORMAL.\hfill$\square$

The supplied script \texttt{normal\_form\_check.py} independently
enumerates all finite orbit classes for $2\leq n\leq7$.  It checks
the count and inverse-transpose action exactly in residue arithmetic;
it is a finite computational check, not a substitute for the proof.
Its fixed counts are $1,0,8,0,192,0$, respectively; for $n=2$,
inverse transpose is inner and the outer automorphism group is trivial.

\subsection{The tangent-rank questions (scan 023)}
Write $\chi_{\mathfrak g}:\mathfrak g\to\mathfrak t/W$ and
$\chi_G:G\to T/W$.  The complete Lie-algebra answer is the
Richardson computation, completed at $E_8(a_2)$ by later published
work.  For $x=s+e$ and $\mathfrak l=\mathfrak g^s$,
Richardson (1987), Proposition~1.2, p.~229, gives
\[
 \operatorname{rank}(d\chi_{\mathfrak g})_x
 =\dim\mathfrak z(\mathfrak l)+
   \sum_j\operatorname{rank}(d\chi_{\mathfrak l_j})_{e_j},
 \qquad [\mathfrak l,\mathfrak l]=\bigoplus_j\mathfrak l_j.
\]
For a nilpotent Jordan partition $\lambda_1\geq\cdots\geq\lambda_k$,
the classical ranks are
\[
\begin{array}{c|c|l}
 A_r&\lambda_1-1&\text{Richardson, Thm.~4.2.1, p.~233}\\
 B_r,C_r&\lfloor\lambda_1/2\rfloor&\text{Thm.~4.3.3, p.~233}\\
 D_r&\lfloor\lambda_1/2\rfloor+
     \mathbf1_{\{k=2,\ \lambda_1\text{ odd}\}}&
     \text{Lemma~4.4.1 and Thm.~4.4.2, p.~234}.
\end{array}
\]
For the exceptional types, the complete finite rank table is
Richardson's Appendix, pp.~239--240, interpreted by Lemma~2.1,
Proposition~2.2 and Corollary~2.3, p.~230, with the precise corrections
in accepted CC-RK02--03: $E_8(a_1)$ has rank $7$, the omitted
$E_6\subset E_8$ class has rank $3$, the final regular $E_8$
exponent is $29$ rather than the printed $27$, and $E_8(a_2)$ has
rank $6$.  The last rank is Dinar (2020), \emph{Journal of Geometry and Physics}
155, article 103690, DOI \texttt{10.1016/j.geomphys.2020.103690};
the inspected author version is arXiv:1905.01829v4,
Corollary~5.5, p.~14 (proof p.~15).  A separate deduction is
CC-RK03 from Johnson--Sommers, Theorem~1.3(1)--(2).
These exact locators, together with the semisimple-centralizer formula,
determine the rank at every Lie-algebra element.  Richardson expressly
calls his article a statement of results with proof outlines (p.~228);
we do not misdescribe it as a full proof reconstructed here.

\paragraph{The subregular speculation is false in every proposed
nonsimply-laced case.}
There is a unique subregular nilpotent orbit in every complex simple
type, and the rank there is $r-1$.  The exact attribution is
Richardson (1987), \S5 opening paragraph and Lemma~5.1, p.~234.
Richardson attributes the rank theorem to Slodowy, his reference~[13],
\emph{Simple Singularities and Simple Algebraic Groups}, LNM~815
(1980).  The nilpotent/unipotent transfer of both uniqueness and rank
is accepted CC-RK04--05; thus the unique subregular unipotent orbit
also has rank $r-1$ in all simple types and all central forms.
This refutes both the strict inequality proposed for unequal root
lengths and the proposed characterization of ADE by uniqueness.
For $A_1$, the subregular orbit is zero and the formula gives zero.

\paragraph{The requested explicit $\mathfrak{sp}_4$ instance
(BR-12302).}
Put
\[
 J=\begin{pmatrix}0&I_2\\-I_2&0\end{pmatrix},\qquad
 N=\begin{pmatrix}0&I_2\\0&0\end{pmatrix},\qquad
 Y=\begin{pmatrix}A&B\\C&-A^{\mathsf T}\end{pmatrix},
 \quad B=B^{\mathsf T},\ C=C^{\mathsf T}.
\]
Then $[Y,N]=\left(\begin{smallmatrix}-C&A+A^{\mathsf T}\\0&C
\end{smallmatrix}\right)$.  Its kernel has dimension $3+1=4$,
so the orbit dimension is $10-4=6=10-2-2$.  With quotient
coordinates $c_2=-\operatorname{Tr}(Y^2)/2$, $c_4=\det Y$,
\[
 dc_2|_N(Y)=-\operatorname{Tr}(NY)=-\operatorname{Tr}C\ne0,
 \qquad dc_4|_N=0,
\]
where the last equality follows because $N$ has rank two and all
$3$ by $3$ cofactors vanish.  Thus the rank is $1=r-1$, not less.
For $U=I+N\in\operatorname{Sp}_4$, the same orbit dimension holds.
On its group quotient use $t=\operatorname{Tr}g$ and
$e=c_2(g)-2$.  At tangent vector $UY$,
\[
 dt_U(UY)=\operatorname{Tr}(NY),\qquad
 de_U(UY)=2\operatorname{Tr}(NY).
\]
The adjoint quotient replaces $(t,e)$ by $(t^2,e)$; since $t(U)=4$
this has invertible differential there.  Rank one therefore also
holds at $[U]\in\operatorname{PSp}_4$, the original adjoint-group
setting.  This is the accepted counterexample, restated explicitly.

\paragraph{The general group quotient needs its own formula.}
To answer the letter away from the unipotent locus, use the finite
character formula BR-12305--08, not a global identification of the
Lie and group quotient maps.  Write $G=\widetilde G/Z$ with
$\widetilde G$ simply connected, choose a lift $\widetilde g$, and
let $\rho_i$, $\chi_i=\operatorname{Tr}\rho_i$, and $\xi_i$
be the fundamental representations, characters, and central
characters.  Define
\[
 H=\{z\in Z:\xi_i(z)\chi_i(\widetilde g)=\chi_i(\widetilde g)
                   \text{ for all }i\},\qquad
 J_H=\{i:\xi_i|_H=1\}.
\]
Then
\[
 \operatorname{rank}(d\chi_G)_g=
 \dim\operatorname{span}\{Y\mapsto
 \operatorname{Tr}(\rho_i(\widetilde g)d\rho_i(Y)):i\in J_H\}.
\]
The independent finite calculation applies in every type and for
arbitrary Jordan decomposition.  Its primary invariant-ring input
is Steinberg (1965), Theorem~6.1(b), p.~62, and (6.9)--(6.10),
pp.~63--64; the proof of the displayed formula is checkpoint~09,
\texttt{research/root\_group\_rank.tex}, BR-12305--06.
In type $\operatorname{SL}_n/\mu_c$, for a lift $A$ with
characteristic coefficients $c_i(A)$, let
\[
 d=\gcd(c,\{i:1\leq i<n,\ c_i(A)\ne0\}).
\]
With the empty-set convention $d=c$, BR-12307 gives
$\operatorname{rank}(d\chi_G)=\deg m_{A^d}-1$.
For the projective class of
$Q=\operatorname{diag}(1,\zeta,\ldots,\zeta^{n-1})$, this is
zero despite regular semisimplicity (BR-12308); its simply connected
lift has rank $n-1$.  The distinction in central form is material.

\subsection{Dimensions and rational singularities (scan 023)}
Let $\mathcal B_x$ be the variety of Borels whose Lie algebra contains
$x$, or of Borels containing $x$ in the group version.  In the complex
setting,
\[
 \dim\mathcal B_x=\frac{\dim Z_G(x)-r}{2}
 =\frac{\dim G-r-\dim(Gx)}2.
\]
Thus an orbit of dimension $2N-2i$, with $2N=\dim G-r$, has
Borel variety of dimension exactly $i$, as the letter predicts.
For group elements this is Steinberg (1976),
\emph{On the desingularization of the unipotent variety},
Theorems~4.6 (p.~217) and 6.5 (p.~223).
For the Lie algebra, split off the semisimple part in its connected
centralizer and apply the nilpotent/unipotent correspondence; the
centralizer reduction is Springer (1984), Lemma~1, p.~273.
This is the accepted dimension result in
\texttt{research/root\_springer\_archive.tex}, equation~(123.1).

The rational-singularity question is affirmative in precisely the
scope of the letter: all fibres of the complex Lie-algebra quotient,
and all fibres of the group quotient when the group is simply
connected.  The exact accepted statement is F-SM-23 in
\texttt{research/source\_cote123\_rational.tex}.  Its Lie-algebra
input is Kraft--Procesi (1982), \S0.3--0.4, p.~540, and its
nilpotent/unipotent input is Hesselink (1976), Theorem~A and \S2,
p.~249.  For a group fibre with semisimple part represented by $s$,
put $L=Z_G(s)$; in the simply connected case it is connected.
The fibre is $G\times^L\mathcal U(L)$, and the induced
$L$-equivariant resolution of $\mathcal U(L)$ proves the required
vanishing by flat base change and descent.  Reducedness and normality
of the group fibre follow from Steinberg (1965),
Theorem~6.11 and \S6.16, pp.~64--65, and Theorem~8.1, p.~70,
as written in F-SM-23.  Hence for a resolution $f:F'\to F$,
$f_*\mathcal O_{F'}=\mathcal O_F$ and
$R^if_*\mathcal O_{F'}=0$ for every $i>0$.  A claim for arbitrary
positive characteristic or arbitrary central form is not added.

\subsection{Where these answers live in Grothendieck's frameworks}
These are statements about algebraic groups and their projective
schemes of Borel subgroups.  In SGA~3, Expos\'e~XXVI,
Proposition~1.2 identifies $G/B$ as a smooth projective scheme;
Theorem~4.3.2 and Definition~4.3.3 define opposition, and
Remark~4.3.6 describes its open Bruhat cell.  Pinnings are
Expos\'e~XXIII, Definition~1.1.  Expos\'e~XXIV, Theorem~1.3(iii)
and Lemma~1.4 give the splitting of algebraic automorphisms into
inner and pinned root-datum automorphisms used here.
The apposition configuration is a finite orbit problem on pairs
of torus-fixed points of the flag scheme $G/B$; opposition is the
open Bruhat orbit.  In type A its defining conditions are explicit
nonvanishing determinants over $\mathbf Z[\zeta_n,1/n]$.
The exclusion verified above is about complex algebraic
outer automorphisms of this configuration.  No automatic extension
of its answer to all fibres over $\operatorname{Spec}\mathbf Z$
is claimed.

The character formula belongs to the representation ring
$R(\widetilde G)$, which is a $\lambda$-ring under exterior powers.
For type A its fundamental characters are exactly
$\chi_{\wedge^i V}=\lambda^i[V]$ evaluated at the group element;
this explains the characteristic-coefficient calculation.  The
Springer fibre statements concern proper algebraic schemes and
their cohomology.  The positive-characteristic, $\ell$-adic purity
and root-of-unity refinements used in ALPOGE\_BRIDGES are separate
published inputs (Springer 1984, Theorem~1, pp.~272--273,
Corollary~1 to Theorem~2 and Lemma~4, p.~280); they are not
consequences of the complex rank calculation alone.  None of the
new normalizer or fixed-class counts proves a new assertion about
motivic $L$-functions or the Tate conjecture.  The singularity
answer has exactly the higher-direct-image formulation printed
on scan 023.

\subsection{ARCH-123-SOURCE: bounded literature audit}
The following is a source-inspection result, not a proof of priority.
Kostant (1959), Theorem~8.6 and Corollary~8.6, were checked in the
OCR of the author's collected-paper reprint; the publisher identifies
the chapter as pp.~130--189, DOI \texttt{10.1007/b94535\_11}.
That OCR omits the pages containing the original \S7.3, so no claim
to have directly read those pages is made.  The supplied complete
Kostant (2010) LaTeX and the author preprint were checked at the
apposition statements, principal-element classification, and
normalizer equation (0.6).  The rank attribution and classical
formulas were read directly in the ANU PDF of Richardson (1987).
Steinberg (1965), \S9, was read at Theorem~9.4, its proof, and
Remark~9.6.  Slodowy's publisher contents and freely available
back matter/index were searched; the body was not available on
that route.  Richardson's explicit attribution supplies the
verified Slodowy credit without an invented internal lemma number.

Additional primary checks covered Guest--Ho,
\emph{Kostant, Steinberg, and the Stokes matrices of the tt*-Toda
equations}, arXiv:1802.01126, especially its apposition and Borel
occurrences; Elashvili--Kac--Vinberg,
\emph{Cyclic elements in semisimple Lie algebras},
arXiv:1205.0515, introduction and \S6; and the Coxeter-torus
terminology in \emph{Orthogonality relations for deep level
Deligne--Lusztig schemes of Coxeter type} (2024),
Definition~2.6.1 and Lemma~2.6.2.  Those works concern related
Coxeter structures, but the inspected passages do not state this
apposition-quadruple opposition table or the outer-invariant
exclusion.  Searches combined ``apposition'', ``Borel'',
``opposite'', ``outer'', ``distinguished'', and ``quadruples''.
They did not locate an exact published antecedent of those two
assertions.  This is explicitly limited search evidence and does
not certify that they are new.  The mathematical answers in the
companion use their supplied proofs independently of that historical
question.
